\documentclass[12pt,a4paper]{report}

% Kodowanie i język
\usepackage[utf8]{inputenc}
\usepackage[T1]{fontenc}
\usepackage[polish]{babel}
\usepackage{lmodern}

% Marginesy
\usepackage[margin=2.5cm]{geometry}

% Grafika i kolory
\usepackage{graphicx}
\usepackage{xcolor}

% Tabele
\usepackage{booktabs}
\usepackage{longtable}
\usepackage{tabularx}
\usepackage{array}
\usepackage{multirow}

% Listy
\usepackage{enumitem}

% Kod źródłowy
\usepackage{listings}
\usepackage{lstautogobble}

% Hiperłącza
\usepackage[hidelinks,colorlinks=false]{hyperref}

% Nagłówki i stopki
\usepackage{fancyhdr}

% Interlinea
\usepackage{setspace}
\onehalfspacing

% Wcięcia
\usepackage{parskip}
\setlength{\parskip}{0.5em}

% Definicja koloru dla kodu
\definecolor{codeblue}{RGB}{0,102,204}
\definecolor{codegray}{RGB}{128,128,128}
\definecolor{codegreen}{RGB}{0,128,0}
\definecolor{codebg}{RGB}{248,248,248}
\definecolor{codeborder}{RGB}{220,220,220}

% Styl dla kodu
\lstdefinestyle{mystyle}{
    backgroundcolor=\color{codebg},
    commentstyle=\color{codegreen},
    keywordstyle=\color{codeblue}\bfseries,
    numberstyle=\tiny\color{codegray},
    stringstyle=\color{codegreen!60!black},
    basicstyle=\ttfamily\small,
    breakatwhitespace=false,
    breaklines=true,
    captionpos=b,
    keepspaces=true,
    numbers=left,
    numbersep=8pt,
    showspaces=false,
    showstringspaces=false,
    showtabs=false,
    tabsize=2,
    frame=single,
    rulecolor=\color{codeborder},
    autogobble=true
}
\lstset{style=mystyle}

% Styl PHP
\lstdefinelanguage{PHP}{
    morekeywords={class,function,public,private,protected,static,return,if,else,foreach,new,use,namespace,extends,implements,interface,abstract,final,try,catch,throw,null,true,false,array,string,int,float,bool,void,mixed},
    sensitive=true,
    morecomment=[l]{//},
    morecomment=[s]{/*}{*/},
    morestring=[b]",
    morestring=[b]',
}

% Nagłówki
\pagestyle{fancy}
\fancyhf{}
\fancyhead[L]{\small Strava Clone -- Sprawozdanie projektowe}
\fancyhead[R]{\small\thepage}
\fancyfoot[C]{}
\renewcommand{\headrulewidth}{0.4pt}

% Głębokość spisu treści
\setcounter{tocdepth}{3}
\setcounter{secnumdepth}{3}


\begin{document}

\begin{titlepage}
    \centering
    \vspace*{2cm}

    {\LARGE\bfseries Sprawozdanie z projektu programistycznego\par}
    \vspace{1.5cm}

    {\Huge\bfseries Strava Clone\par}
    \vspace{0.5cm}
    {\large Aplikacja do śledzenia aktywności fizycznych\par}

    \vspace{1cm}
    {\large\bfseries Projektowanie i Programowanie Systemów Internetowych I\par}

    \vspace{1.5cm}

    \begin{tabular}{ll}
        \textbf{Technologia backendu:}  & Laravel 12 (PHP 8.2+)       \\[4pt]
        \textbf{Technologia frontendu:} & Vue 3 + TypeScript           \\[4pt]
        \textbf{Baza danych:}           & PostgreSQL 18                \\[4pt]
        \textbf{Repozytorium:}          & github.com/Console-dungeon/ \\
                                        & strava-clone                 \\
    \end{tabular}

    \vfill

    {\large \today\par}
\end{titlepage}

\tableofcontents
\thispagestyle{empty}
\cleardoublepage

\chapter{Opis funkcjonalny systemu}
\label{ch:funkcjonalny}

\section{Cel i zakres projektu}

Strava Clone jest aplikacją webową przeznaczoną do rejestrowania, analizowania i wizualizacji aktywności fizycznych użytkownika. Stanowi odwzorowanie kluczowych funkcjonalności popularnej platformy Strava, z naciskiem na integrację z urządzeniami Garmin oraz wsparciem dla plików GPS.

Aplikacja adresowana jest do osób aktywnych fizycznie -- biegaczy, rowerzystów i pływaków -- którzy chcą mieć wgląd w historię swoich treningów i monitorować postępy w jednym miejscu.

\section{Moduły funkcjonalne}

\subsection{Rejestracja i uwierzytelnianie}

System obsługuje pełny cykl życia konta użytkownika:
\begin{itemize}[noitemsep]
    \item rejestracja nowego konta z walidacją adresu e-mail,
    \item logowanie i wylogowanie z obsługą sesji,
    \item reset hasła przez e-mail (token jednorazowy),
    \item weryfikacja adresu e-mail,
    \item zmiana hasła i danych profilowych,
    \item usunięcie konta z usunięciem wszystkich powiązanych danych.
\end{itemize}

Moduł oparty jest na pakiecie \textit{Laravel Breeze} i nie wymaga modyfikacji.

\subsection{Zarządzanie aktywnościami}

Rdzeń funkcjonalności aplikacji umożliwia:
\begin{itemize}[noitemsep]
    \item ręczne dodanie aktywności (typ, dystans, czas, notatki),
    \item import aktywności z pliku GPX -- aplikacja parsuje trasę i oblicza parametry,
    \item przeglądanie listy aktywności z filtrowaniem po typie i paginacją,
    \item wyświetlanie szczegółów aktywności wraz z interaktywną mapą trasy,
    \item usuwanie aktywności.
\end{itemize}

Obsługiwane typy aktywności: bieg (\textit{running}), jazda na rowerze (\textit{cycling}), pływanie (\textit{swimming}) oraz inne.

\subsection{Integracja z Garmin Connect}

Aplikacja umożliwia:
\begin{itemize}[noitemsep]
    \item jednorazowy pełny import historii aktywności z konta Garmin Connect (do 1000 aktywności),
    \item synchronizację nowych aktywności (brakujące od ostatniego importu),
    \item pobranie szczegółowych danych aktywności Garmin: średnie i maksymalne tętno, spalone kalorie, prędkość maksymalna,
    \item przechowywanie danych uwierzytelniających Garmin zaszyfrowanych w bazie danych.
\end{itemize}

\subsection{Dashboard i statystyki}

Strona główna zalogowanego użytkownika prezentuje:
\begin{itemize}[noitemsep]
    \item karty podsumowania: łączny dystans, łączny czas, średnia prędkość,
    \item wykres liniowy aktywności w czasie (za pomocą biblioteki Unovis),
    \item tabelę ostatnich aktywności z najważniejszymi parametrami.
\end{itemize}

\subsection{Wizualizacja trasy na mapie}

Dla aktywności zaimportowanych z pliku GPX lub Garmina dostępna jest interaktywna mapa oparta na bibliotece Leaflet z OpenStreetMap, prezentująca:
\begin{itemize}[noitemsep]
    \item narysowaną trasę aktywności,
    \item marker punktu startowego i końcowego trasy.
\end{itemize}

\subsection{Logi aktywności użytkownika}

System rejestruje kluczowe akcje użytkownika w tabeli \texttt{activity\_logs}:
\begin{itemize}[noitemsep]
    \item uruchomienie synchronizacji Garmin,
    \item zakończenie synchronizacji Garmin,
    \item utworzenie aktywności,
    \item usunięcie aktywności.
\end{itemize}

\subsection{Internacjonalizacja}

Interfejs dostępny jest w dwóch językach:
\begin{itemize}[noitemsep]
    \item \textbf{polskim} -- domyślnym (zmienna \texttt{APP\_LOCALE=pl}),
    \item \textbf{angielskim} -- jako język zapasowy.
\end{itemize}
Tłumaczenia przechowywane są w plikach TypeScript \texttt{pl.ts} i \texttt{en.ts}.

\subsection{Profil użytkownika}

Moduł profilu pozwala użytkownikowi:
\begin{itemize}[noitemsep]
    \item zaktualizować imię i adres e-mail,
    \item zmienić hasło,
    \item podłączyć i zapisać dane logowania do konta Garmin Connect,
    \item trwale usunąć konto.
\end{itemize}

\section{Diagram przepływu użytkownika}

\begin{enumerate}
    \item Nowy użytkownik rejestruje konto i weryfikuje e-mail.
    \item Po zalogowaniu widzi pusty dashboard z zaproszeniem do dodania aktywności.
    \item Użytkownik może ręcznie wpisać aktywność, wgrać plik GPX lub powiązać konto Garmin.
    \item Po imporcie aktywności dashboard wyświetla statystyki i wykresy.
    \item Kliknięcie aktywności otwiera szczegóły z mapą trasy (jeśli dostępna).
    \item Regularne synchronizacje Garmin mogą być uruchamiane jednym kliknięciem.
\end{enumerate}

% ============================================================

\chapter{Opis technologiczny}
\label{ch:technologiczny}

\section{Architektura ogólna}

Aplikacja zbudowana jest w architekturze monolitycznej z wbudowanym podziałem na warstwy, opartej na wzorcu MVC rozszerzonym o warstwę serwisów. Komunikacja między backendem a frontendem odbywa się przez protokół \textbf{Inertia.js}, który eliminuje potrzebę budowania osobnego REST API -- backend zwraca komponenty Vue zamiast JSON.

\begin{center}
\begin{tabular}{|l|l|}
\hline
\textbf{Warstwa}       & \textbf{Technologia}               \\
\hline
Backend framework      & Laravel 12 (PHP 8.2+)              \\
Frontend framework     & Vue 3 + TypeScript                 \\
Protokół SPA           & Inertia.js 2.0                     \\
Baza danych            & PostgreSQL 18                      \\
Bundler                & Vite 7                             \\
Serwer deweloperski    & Laravel Sail (Docker)              \\
\hline
\end{tabular}
\end{center}

\section{Backend}

\subsection{Laravel 12}

Framework PHP w wersji 12 dostarcza:
\begin{itemize}[noitemsep]
    \item \textbf{Routing} -- deklaratywny w \texttt{routes/web.php} i \texttt{routes/auth.php},
    \item \textbf{Eloquent ORM} -- mapowanie obiektowo-relacyjne dla PostgreSQL,
    \item \textbf{Form Requests} -- walidacja danych wejściowych w osobnych klasach,
    \item \textbf{Queue} -- asynchroniczne przetwarzanie zadań (driver: \texttt{database}),
    \item \textbf{Cache} -- buforowanie danych (driver: \texttt{database}),
    \item \textbf{Szyfrowanie} -- kolumna \texttt{garmin\_password} przechowywana zaszyfrowana (\texttt{encrypted} cast).
\end{itemize}

\subsection{Wzorzec Controller--Service}

Kontrolery są celowo cienkie: nie zawierają logiki biznesowej, jedynie delegują pracę do warstwy serwisów.

\begin{lstlisting}[language=PHP, caption=Przykładowy kontroler delegujący do serwisu]
class DashboardController extends Controller
{
    public function __construct(
        private readonly DashboardService $dashboard
    ) {}

    public function __invoke(Request $request): Response
    {
        return Inertia::render('Dashboard', [
            'data' => $this->dashboard->getData(
                auth()->guard()->user()
            ),
        ]);
    }
}
\end{lstlisting}

\subsection{Inertia.js}

Inertia.js pełni rolę warstwy pośredniczącej między backendem a frontendem. Backend nie zwraca JSON, lecz wywołuje \texttt{Inertia::render('NazwaStrony', \$props)}, co powoduje załadowanie odpowiedniego komponentu Vue z danymi jako \textit{props}. Nawigacja po aplikacji odbywa się bez przeładowania strony (SPA), przy czym serwer zawsze kontroluje dane i routowanie.

\subsection{Kluczowe serwisy}

\subsubsection{GarminService}

Odpowiada za komunikację z Garmin Connect przez pakiet \texttt{alexoid/laravel-garmin}:
\begin{itemize}[noitemsep]
    \item \texttt{importActivities()} -- pobiera do 1000 aktywności i zapisuje je w bazie,
    \item \texttt{syncActivities()} -- pobiera wyłącznie aktywności nowsze niż ostatnia zapisana.
\end{itemize}
Prędkości z Garmin (w m/s) są przeliczane na km/h przez mnożnik $\times 3{,}6$.

\subsubsection{GpxParser}

Parsuje pliki GPX (format XML) i wyodrębnia:
\begin{itemize}[noitemsep]
    \item listę punktów trasy (szerokość/długość geograficzna, wysokość, czas),
    \item całkowity dystans -- obliczony wzorem Haversine'a,
    \item przewyższenie dodatnie (\textit{elevation gain}),
    \item czas aktywności, datę, typ.
\end{itemize}
Dla wydajności trasa jest upraszczana do maksymalnie 500 punktów.

\subsubsection{DashboardService}

Agreguje dane dla dashboardu użytkownika:
\begin{itemize}[noitemsep]
    \item łączny dystans, łączny czas, średnia prędkość,
    \item lista ostatnich aktywności,
    \item dane szeregów czasowych dla wykresu.
\end{itemize}

\subsubsection{ActivityLogger}

Rejestruje akcje użytkownika w tabeli \texttt{activity\_logs}. Zdefiniowane akcje:
\texttt{GARMIN\_SYNC\_STARTED}, \texttt{GARMIN\_SYNC\_COMPLETED}, \texttt{ACTIVITY\_CREATED}, \texttt{ACTIVITY\_DELETED}.

\subsection{Uwierzytelnianie}

System uwierzytelniania oparty na \textbf{Laravel Breeze}:
\begin{itemize}[noitemsep]
    \item sesje przechowywane w bazie danych (tabela \texttt{sessions}),
    \item middleware \texttt{auth} chroni wszystkie trasy wymagające zalogowania,
    \item middleware \texttt{HandleInertiaRequests} eksponuje dane zalogowanego użytkownika (\texttt{auth.user}) do każdej strony Vue jako \textit{shared props}.
\end{itemize}

\section{Frontend}

\subsection{Vue 3 + TypeScript}

Interfejs użytkownika napisany jest w Vue 3 z Composition API (\texttt{<script setup lang="ts">}). TypeScript działa w trybie \textit{strict} -- wszystkie zmienne są typowane, użycie \texttt{any} jest zakazane.

\subsection{Kluczowe biblioteki frontendowe}

\begin{center}
\begin{longtable}{|l|p{8cm}|}
\hline
\textbf{Biblioteka}       & \textbf{Zastosowanie}                              \\
\hline
Reka UI                   & Beznakładowe prymitywy UI (Dialog, DropdownMenu, Table i in.) \\
\hline
Tailwind CSS v4           & Utility-first CSS -- stylowanie komponentów       \\
\hline
Lucide Vue Next           & Ikony wektorowe SVG                               \\
\hline
Vue i18n 11               & Internacjonalizacja -- klucze tłumaczeń w szablonach \\
\hline
Vue Sonner                & Powiadomienia toast                               \\
\hline
Unovis                    & Wykresy liniowe (biblioteka wizualizacji danych)  \\
\hline
Leaflet + OpenStreetMap   & Interaktywna mapa tras GPS                        \\
\hline
TanStack Vue Table        & Tabela danych z sortowaniem i filtrowaniem        \\
\hline
Ziggy                     & Generowanie URL z nazwanych tras Laravel w JS     \\
\hline
VueUse                    & Utilities Composition API (hooks pomocnicze)      \\
\hline
\end{longtable}
\end{center}

\subsection{Struktura plików frontendowych}

\begin{lstlisting}[caption=Struktura katalogu resources/js/]
resources/js/
├── Pages/           # Komponenty stron (1:1 z routami)
│   ├── Activities/  # Widoki aktywności
│   ├── Auth/        # Strony logowania/rejestracji
│   ├── Dashboard/   # Komponenty dashboardu
│   └── Profile/     # Profil użytkownika
├── Components/
│   └── ui/          # Prymitywy Reka UI (nie edytować)
├── Layouts/         # AppLayout, GuestLayout
├── i18n/            # pl.ts, en.ts, index.ts
└── app.ts           # Punkt wejściowy aplikacji
\end{lstlisting}

\subsection{Konwencje frontendowe}

\begin{itemize}[noitemsep]
    \item Tylko \texttt{<script setup lang="ts">} -- zakaz używania Options API,
    \item alias ścieżki \texttt{@/} mapuje na \texttt{resources/js/},
    \item nawigacja przez komponent \texttt{<Link>} z Inertia (nie tagi \texttt{<a>}),
    \item tryb ciemny (dark mode) -- każdy nowy komponent musi wspierać oba warianty,
    \item formatowanie przez Prettier (pojedyncze cudzysłowy, średniki, 80 znaków, 2 spacje wcięcia).
\end{itemize}

\section{Środowisko deweloperskie}

\subsection{Docker / Laravel Sail}

Środowisko lokalne opiera się na \textbf{Laravel Sail} -- zestawie skryptów Docker skonfigurowanych dla Laravel. Sail uruchamia kontenery:
\begin{itemize}[noitemsep]
    \item \textbf{PHP/Laravel} -- serwer aplikacji,
    \item \textbf{PostgreSQL 18} -- baza danych,
    \item \textbf{Adminer} -- graficzny interfejs bazy danych dostępny pod adresem \texttt{http://localhost:8080},
    \item \textbf{Node.js} -- kompilacja frontendu przez Vite.
\end{itemize}

\subsection{Dev Container}

Projekt skonfigurowany jest dla \textbf{VS Code Dev Containers} -- jednorazowe otwarcie projektu w kontenerze zapewnia gotowe do pracy środowisko bez ręcznej instalacji zależności.

\subsection{Haki pre-commit}

\textbf{Husky} uruchamia automatycznie \texttt{pnpm lint-staged} przed każdym commitem:
\begin{itemize}[noitemsep]
    \item ESLint z auto-naprawą dla plików \texttt{.ts} / \texttt{.vue},
    \item Prettier dla formatowania kodu,
    \item Laravel Pint dla PHP.
\end{itemize}
Pomijanie haków przez \texttt{--no-verify} jest zakazane.

% ============================================================

\chapter{Wdrożone zagadnienia kwalifikacyjne}
\label{ch:zagadnienia}

Poniżej wyszczególnione zostały zagadnienia techniczne i architektoniczne wdrożone w projekcie, stanowiące o dojrzałości technicznej rozwiązania.

\section{Architektura i wzorce projektowe}

\begin{enumerate}

\item \textbf{Wzorzec MVC z warstwą serwisów} \\
Aplikacja stosuje wzorzec Model--Widok--Kontroler rozszerzony o dedykowaną warstwę serwisów. Kontrolery delegują całą logikę biznesową do klas serwisowych (\texttt{GarminService}, \texttt{DashboardService}, \texttt{GpxParser}), pozostając cienkimi adapterami między routingiem a logiką.

\item \textbf{Architektura SPA z Inertia.js} \\
Zamiast tradycyjnego REST API + osobna aplikacja SPA, projekt stosuje Inertia.js jako protokół łączący Laravel z Vue 3. Backend renderuje komponenty Vue z danymi bez konieczności budowania i utrzymywania API JSON. Nawigacja odbywa się bez przeładowania strony.

\item \textbf{Dependency Injection} \\
Serwisy wstrzykiwane są przez konstruktory kontrolerów za pośrednictwem kontenera IoC Laravel. Przykład: \texttt{DashboardController} deklaruje zależność od \texttt{DashboardService} przez typowanie konstruktora.

\item \textbf{Repozytorium zdarzeń -- Activity Logs} \\
Zaimplementowano niezmienny rejestr zdarzeń (append-only log) w tabeli \texttt{activity\_logs}, rejestrujący kluczowe akcje z sygnaturą czasową. Tabela nie posiada kolumny \texttt{updated\_at} i nigdy nie jest aktualizowana -- jedynie dopisywana.

\end{enumerate}

\section{Bezpieczeństwo}

\begin{enumerate}[resume]

\item \textbf{Szyfrowanie wrażliwych danych} \\
Hasła do konta Garmin Connect przechowywane są w bazie danych z użyciem wbudowanego mechanizmu szyfrowania Laravel (cast \texttt{encrypted}). Klucz szyfrowania pochodzi z \texttt{APP\_KEY}.

\item \textbf{Haszowanie haseł} \\
Hasła użytkowników haszowane są algorytmem bcrypt (12 rund, konfigurowane przez \texttt{BCRYPT\_ROUNDS}), nigdy nie przechowywane w postaci jawnej.

\item \textbf{Izolacja danych między użytkownikami} \\
Każde zapytanie do tabeli \texttt{activities} scopowane jest do zalogowanego użytkownika (\texttt{WHERE user\_id = ?}). Brak możliwości dostępu do cudzych aktywności.

\item \textbf{Walidacja danych wejściowych} \\
Dane formularzy walidowane są w dedykowanych klasach \texttt{Http/Requests/} przed dotarciem do kontrolerów. Walidacja po stronie serwera jest obowiązkowa -- nie polegamy wyłącznie na walidacji frontendu.

\item \textbf{Autoryzacja tras} \\
Middleware \texttt{auth} chroni wszystkie trasy wymagające zalogowania. Trasy dostępne dla gości (rejestracja, logowanie) otoczone są middleware \texttt{guest}, uniemożliwiającym dostęp po zalogowaniu.

\item \textbf{CSRF Protection} \\
Laravel automatycznie weryfikuje tokeny CSRF dla wszystkich żądań zmieniających stan (POST, PUT, PATCH, DELETE). Inertia.js dołącza token do nagłówków żądań.

\end{enumerate}

\section{Integracje zewnętrzne}

\begin{enumerate}[resume]

\item \textbf{Integracja z Garmin Connect API} \\
Wdrożono dwutrybow\`a integrację z platformą Garmin Connect przez pakiet \texttt{alexoid/laravel-garmin}: pełny import historii oraz inkrementalna synchronizacja nowych aktywności. Przeliczanie jednostek (m/s $\to$ km/h).

\item \textbf{Parser plików GPX} \\
Zaimplementowano własny parser plików GPS Exchange Format (XML). Parser oblicza dystans wzorem Haversine'a, wyznacza przewyższenie dodatnie i upraszcza trasę algorytmem Douglas--Peucker do 500 punktów dla wydajności.

\item \textbf{Mapa interaktywna -- Leaflet + OpenStreetMap} \\
Trasy aktywności wizualizowane są na interaktywnej mapie z kafli OpenStreetMap, bez konieczności korzystania z płatnych API map.

\end{enumerate}

\section{Frontend i UX}

\begin{enumerate}[resume]

\item \textbf{Tryb ciemny (Dark Mode)} \\
Interfejs w pełni wspiera tryb ciemny Tailwind CSS. Każdy komponent definiuje warianty dla jasnego i ciemnego motywu.

\item \textbf{Internacjonalizacja (i18n)} \\
Aplikacja obsługuje dwa języki (polski i angielski) z użyciem Vue i18n. Żadny widoczny ciąg znaków nie jest hardkodowany w szablonach -- wszystko przechodzi przez klucze tłumaczeń.

\item \textbf{Powiadomienia Toast} \\
System flash-wiadomości zintegrowany między backendem (Inertia flash) a frontendem (Vue Sonner). Backend ustawia komunikat, frontend wyświetla powiadomienie bez dodatkowego kodu po stronie stron.

\item \textbf{Komponenty headless UI -- Reka UI} \\
Prymitywy interfejsu (okna dialogowe, menu rozwijane, tabele) zbudowane są na bazie biblioteki Reka UI (fork Radix Vue), zapewniającej pełną dostępność (a11y) i poprawne zarządzanie focusem klawiatury.

\item \textbf{Wykresy danych -- Unovis} \\
Dashboard prezentuje wykres liniowy aktywności z użyciem biblioteki Unovis, dedykowanej do wizualizacji danych w aplikacjach Vue.

\end{enumerate}

\section{Jakość kodu i DevOps}

\begin{enumerate}[resume]

\item \textbf{TypeScript Strict Mode} \\
Cały frontend pisany jest w TypeScript z włączonym trybem \textit{strict} (\texttt{vue-tsc} jako sprawdzanie typów). Użycie \texttt{any} jest zakazane.

\item \textbf{Automatyczne formatowanie i lintowanie} \\
Pipeline jakości kodu: Prettier (formatowanie), ESLint (lintowanie TypeScript/Vue), Laravel Pint (formatowanie PHP). Haki Husky uruchamiają \texttt{lint-staged} automatycznie przed każdym commitem.

\item \textbf{Konteneryzacja -- Docker / Laravel Sail} \\
Środowisko deweloperskie w pełni skonteneryzowane, uruchamiane jedną komendą. VS Code Dev Containers pozwala na natychmiastowe uruchomienie projektu bez instalacji zależności na hoście.

\item \textbf{Migracje bazy danych} \\
Każda zmiana schematu bazy to osobny, wersjonowany plik migracji Laravel. Historia zmian schematu jest częścią historii Git i pozwala na odtworzenie dowolnego stanu bazy.

\item \textbf{Graficzny interfejs bazy danych -- Adminer} \\
Do środowiska deweloperskiego dołączony jest kontener \textbf{Adminer} -- webowy klient PostgreSQL dostępny pod adresem \texttt{http://localhost:8080}. Narzędzie umożliwia przeglądanie i edycję danych, podgląd struktury tabel oraz wykonywanie zapytań SQL bez konieczności korzystania z wiersza poleceń. Kontener skonfigurowany jest z domyślnym serwerem \texttt{pgsql} i ciemnym motywem wizualnym.

\item \textbf{Kolejkowanie zadań} \\
Synchronizacja z Garmin wykonywana jest za pośrednictwem kolejki zadań Laravel (driver: \texttt{database}), co pozwala na asynchroniczne wykonywanie długotrwałych operacji bez blokowania odpowiedzi HTTP.

\item \textbf{Leniwy zapis danych trasy GPS (lazy persistence)} \\
Dane trasy GPS dla aktywności Garmin pobierane są z API Garmin tylko przy pierwszym wyświetleniu mapy, a następnie trwale zapisywane w tabeli \texttt{activity\_gpx}. Każde kolejne otwarcie mapy korzysta z bazy danych -- bez ponownego zapytania do API Garmin. Tokeny OAuth Garmin przechowywane są osobno w cache Laravel (\texttt{LaravelCacheTokenStore}).

\end{enumerate}


\chapter{Instrukcja uruchomienia systemu}
\label{ch:uruchomienie}

\section{Wymagania wstępne}

\subsection{Uruchomienie lokalne (Dev Container)}

\begin{itemize}[noitemsep]
    \item \textbf{Docker Desktop} -- wersja 4.0 lub nowsza,
    \item \textbf{Visual Studio Code} z rozszerzeniem \textit{Dev Containers} (\texttt{ms-vscode-remote.remote-containers}),
    \item konto Garmin Connect (opcjonalnie, wymagane do integracji Garmin).
\end{itemize}

\subsection{Uruchomienie zdalne / produkcyjne}

\begin{itemize}[noitemsep]
    \item serwer Linux z PHP 8.2+, Node.js 20+, Composer, PostgreSQL 18,
    \item lub środowisko Docker z możliwością uruchomienia \texttt{docker-compose}.
\end{itemize}

\section{Uruchomienie lokalne (Dev Container -- zalecane)}

\subsection{Krok 1: Sklonowanie repozytorium}

\begin{lstlisting}[language=bash, numbers=none]
git clone https://github.com/Console-dungeon/strava-clone.git
cd strava-clone
\end{lstlisting}

\subsection{Krok 2: Inicjalizacja środowiska}

\begin{lstlisting}[language=bash, numbers=none]
make init
\end{lstlisting}

Komenda \texttt{make init} wykonuje:
\begin{enumerate}[noitemsep]
    \item kopiuje \texttt{.env.example} $\to$ \texttt{.env},
    \item instaluje zależności PHP przez Composer (\texttt{composer install}).
\end{enumerate}

\subsection{Krok 3: Otwarcie w kontenerze}

W VS Code: \texttt{Ctrl+Shift+P} $\to$ \textit{Dev Containers: Reopen in Container}.

VS Code automatycznie:
\begin{enumerate}[noitemsep]
    \item zbuduje kontenery Docker (PHP, PostgreSQL, Node.js),
    \item zainstaluje zależności Node.js,
    \item skonfiguruje środowisko.
\end{enumerate}

\subsection{Krok 4: Migracja bazy danych}

\begin{lstlisting}[language=bash, numbers=none]
php artisan migrate
\end{lstlisting}

\subsection{Krok 5: Uruchomienie serwera deweloperskiego}

\begin{lstlisting}[language=bash, numbers=none]
composer dev
\end{lstlisting}

Komenda uruchamia równolegle (za pomocą \texttt{concurrently}):
\begin{itemize}[noitemsep]
    \item \texttt{php artisan serve} -- serwer PHP na porcie 8000,
    \item \texttt{php artisan queue:listen} -- worker kolejki,
    \item \texttt{pnpm run dev} -- serwer Vite (hot-reload frontendu).
\end{itemize}

Aplikacja dostępna pod adresem: \texttt{http://localhost:8000}

\subsection{Alternatywne uruchomienie przez F5}

Projekt skonfigurowany jest dla debuggera VS Code. Naciśnięcie klawisza \textbf{F5} uruchamia serwer deweloperski z aktywnym debuggerem PHP.

\section{Konfiguracja integracji Garmin (opcjonalne)}

\begin{enumerate}
    \item Zaloguj się do aplikacji i przejdź do strony \textbf{Profil}.
    \item W sekcji \textit{Garmin Connect} wpisz adres e-mail i hasło konta Garmin.
    \item Kliknij \textit{Zapisz} -- dane zostaną zaszyfrowane i zapisane w bazie.
    \item Na dashboardzie użyj przycisku \textbf{Importuj z Garmin} (jednorazowy pełny import) lub \textbf{Synchronizuj} (tylko nowe aktywności).
\end{enumerate}

\section{Uruchomienie zdalne / produkcyjne}

\subsection{Krok 1: Przygotowanie serwera}

Upewnij się, że na serwerze zainstalowane są:
\begin{lstlisting}[language=bash, numbers=none]
php --version   # >= 8.2
node --version  # >= 20
composer --version
psql --version  # PostgreSQL >= 15
\end{lstlisting}

\subsection{Krok 2: Sklonowanie i konfiguracja}

\begin{lstlisting}[language=bash, numbers=none]
git clone https://github.com/Console-dungeon/strava-clone.git
cd strava-clone
cp .env.example .env
\end{lstlisting}

\subsection{Krok 3: Konfiguracja zmiennych środowiskowych}

Edytuj plik \texttt{.env} i ustaw:

\begin{lstlisting}[numbers=none]
APP_ENV=production
APP_DEBUG=false
APP_URL=https://twoja-domena.pl

DB_CONNECTION=pgsql
DB_HOST=localhost
DB_PORT=5432
DB_DATABASE=strava_clone
DB_USERNAME=uzytkownik
DB_PASSWORD=twoje_haslo

QUEUE_CONNECTION=database
SESSION_DRIVER=database
CACHE_STORE=database
\end{lstlisting}

\subsection{Krok 4: Instalacja zależności i build}

\begin{lstlisting}[language=bash, numbers=none]
# Zależności PHP (bez deweloperskich)
composer install --no-dev --optimize-autoloader

# Wygenerowanie klucza aplikacji
php artisan key:generate

# Zależności Node.js i build frontendu
npm install
npm run build

# Migracja bazy danych
php artisan migrate --force

# Optymalizacja konfiguracji
php artisan config:cache
php artisan route:cache
php artisan view:cache
\end{lstlisting}

\subsection{Krok 5: Konfiguracja serwera WWW}

\textbf{Nginx} (przykładowa konfiguracja):

\begin{lstlisting}[numbers=none]
server {
    listen 80;
    server_name twoja-domena.pl;
    root /var/www/strava-clone/public;
    index index.php;

    location / {
        try_files $uri $uri/ /index.php?$query_string;
    }

    location ~ \.php$ {
        fastcgi_pass unix:/run/php/php8.2-fpm.sock;
        fastcgi_param SCRIPT_FILENAME $realpath_root$fastcgi_script_name;
        include fastcgi_params;
    }
}
\end{lstlisting}

\subsection{Krok 6: Uruchomienie workera kolejki}

Worker kolejki jest wymagany do działania synchronizacji Garmin:

\begin{lstlisting}[language=bash, numbers=none]
# Uruchomienie przez Supervisor (zalecane na produkcji)
php artisan queue:work --sleep=3 --tries=3

# Lub jako proces systemd / supervisor
\end{lstlisting}

\subsection{Krok 7: Uprawnienia katalogów}

\begin{lstlisting}[language=bash, numbers=none]
chmod -R 775 storage bootstrap/cache
chown -R www-data:www-data storage bootstrap/cache
\end{lstlisting}

\section{Przydatne komendy}

\begin{center}
\begin{tabular}{|l|l|}
\hline
\textbf{Komenda}                      & \textbf{Działanie}                       \\
\hline
\texttt{composer dev}                 & Uruchomienie serwera deweloperskiego      \\
\texttt{php artisan migrate}          & Wykonanie migracji bazy danych            \\
\texttt{php artisan migrate:rollback} & Cofnięcie ostatniej migracji              \\
\texttt{pnpm lint}                    & Lintowanie kodu TypeScript/Vue            \\
\texttt{pnpm format}                  & Formatowanie kodu Prettier                \\
\texttt{./vendor/bin/pint}            & Formatowanie kodu PHP                     \\
\texttt{pnpm build}                   & Build produkcyjny frontendu               \\
\texttt{php artisan queue:listen}     & Nasłuchiwanie na zadania kolejki          \\
\texttt{php artisan tinker}           & Interaktywna konsola Laravel              \\
\hline
\end{tabular}
\end{center}

% ============================================================

\chapter{Wnioski projektowe}
\label{ch:wnioski}

\section{Podsumowanie realizacji projektu}

Projekt zakończył się zrealizowaniem kompletnej aplikacji webowej do śledzenia aktywności fizycznych. W toku pracy zespołowej zdobyto szereg doświadczeń i obserwacji dotyczących organizacji pracy, doboru narzędzi i metodyki wytwarzania oprogramowania. Poniżej przedstawiono wnioski, które mogą być wartościowe przy realizacji kolejnych projektów podobnej skali.

\section{Kluczowe wnioski}

\subsection{Bezpośrednia komunikacja zespołu jako fundament projektu}

Jednym z najistotniejszych czynników sukcesu projektu okazały się regularne spotkania członków zespołu -- zarówno osobiste, jak i przez wideo. Omawianie bieżących kwestii w czasie rzeczywistym, nawet w formie krótkich sesji online, pozwalało na szybkie uzgodnienie kierunku implementacji, rozwiązywanie nieporozumień co do architektury i podziału zadań oraz uniknięcie sytuacji, w których dwóch programistów pracuje nad tym samym problemem lub -- co gorsza -- w sprzecznych kierunkach.

Próby koordynacji wyłącznie przez komentarze w kodzie, wiadomości tekstowe czy opis zadań w systemie kontroli wersji okazały się niewystarczające przy bardziej złożonych decyzjach projektowych. Bezpośrednia komunikacja, choćby krótka, pozostaje niezastąpiona i praktycznie niezbędna przy omawianiu kluczowych kwestii projektu i wdrażanych rozwiązań.

\subsection{Iteracyjna praca z małymi przyrostami}

Projekt prowadzony był metodą iteracyjną -- zamiast próby zaimplementowania wszystkich funkcjonalności naraz, pracowano nad małymi, samodzielnymi przyrostami. Każda iteracja kończyła się działającą funkcją dostępną w głównej gałęzi projektu. Takie podejście przyniosło kilka wymiernych korzyści:
\begin{itemize}
    \item Błędy pojawiały się w kontekście świeżo dodanego, małego fragmentu kodu -- znacznie łatwiej było zidentyfikować ich źródło i naprawić je, zanim zaciemniły większy fragment systemu.
    \item Historia Git pozostawała czytelna -- każdy commit odpowiadał konkretnej, skończonej funkcji.
    \item Projekt był stale w stanie działającym, co eliminowało ryzyko „dużego wybuchu" na etapie integracji.
\end{itemize}
Regularne, małe wdrożenia minimalizowały ryzyko i ułatwiały naprawę błędów. Podejście monolityczne (duże zmiany rzadko) wielokrotnie prowadziłoby do trudnych do debugowania konfliktów.

\subsection{Praca zespołowa jako akcelerator skali i jakości}

Realizacja projektu w zespole pozwoliła na zbudowanie aplikacji o znacznie szerszym zakresie funkcjonalnym, niż byłoby to możliwe w tym samym czasie w pracy indywidualnej. Podział pracy według kompetencji i zainteresowań -- jeden programista skupia się na backendzie i integracji Garmin, inny na fronendzie i wizualizacji -- pozwolił na równoległe postępy w różnych obszarach systemu. Wzajemna weryfikacja kodu i wspólne rozwiązywanie problemów prowadziły do powstawania rozwiązań wyższej jakości niż praca w izolacji.

Praca zespołowa wymaga inwestycji w komunikację i koordynację, jednak zwrot z tej inwestycji jest wyraźny zarówno w tempie powstawania nowych funkcjonalności, jak i w końcowej jakości produktu.

\subsection{Agenci AI jako katalizator produktywności}

Projekt korzystał z narzędzi opartych na agentach sztucznej inteligencji (m.in. Claude Code) do wspomagania implementacji. Zastosowanie agentów AI przyniosło zauważalne skrócenie czasu potrzebnego na tworzenie nowych funkcjonalności -- szkielety kodu, serwisy, komponenty Vue czy migracje bazy danych powstawały w ułamku czasu, jaki zajęłoby ich ręczne napisanie od zera.

Agenci sprawdzali się szczególnie dobrze przy:
\begin{itemize}[noitemsep]
    \item generowaniu powtarzalnych struktur (migracje, kontrolery, komponenty),
    \item eksploracji nieznanych bibliotek i ich API,
    \item szybkim prototypowaniu rozwiązań do weryfikacji koncepcji,
    \item refaktoryzacji i poprawie jakości istniejącego kodu.
\end{itemize}

Kluczowym wnioskiem jest jednak to, że agenci AI są narzędziem amplifikującym -- zwiększają produktywność doświadczonego programisty, który rozumie generowany kod i jest w stanie go zweryfikować, a nie zastępują samodzielnego myślenia projektowego.

\section{Rekomendacje na przyszłość}

Na podstawie doświadczeń zdobytych podczas realizacji projektu sformułowano następujące rekomendacje dla przyszłych przedsięwzięć:

\begin{enumerate}
    \item \textbf{Ustalenie kanałów komunikacji na początku projektu} -- jasno określić, kiedy używamy wiadomości tekstowej, a kiedy planujemy spotkanie.
    \item \textbf{Zdefiniowanie standardów kodu przed pierwszym commitem} -- konfiguracja linterów, formaterów i haków pre-commit od pierwszego dnia eliminuje późniejsze konflikty stylistyczne.
    \item \textbf{Małe, częste commity zamiast rzadkich dużych paczek zmian} -- ułatwia debugowanie, code review i cofanie błędnych decyzji.
    \item \textbf{Integracja narzędzi AI od początku projektu} -- wcześniejsze oswojenie się z narzędziami AI w kontekście projektu przynosi większe korzyści niż sięganie po nie dopiero przy trudniejszych zadaniach.
\end{enumerate}

\section{Zakończenie}

Projekt pozwolił zbudować kompletną aplikację do śledzenia aktywności fizycznych z integracją zewnętrznych platform sportowych, nowoczesnym frontendem i przemyślaną architekturą backendową. Sukces projektu jest wynikiem połączenia regularnej komunikacji w zespole, iteracyjnego podejścia do wytwarzania oprogramowania, efektywnej współpracy i umiejętnego wykorzystania nowoczesnych narzędzi wspomagania programowania. Doświadczenia zdobyte podczas realizacji projektu stanowią solidną podstawę dla kolejnych, bardziej ambitnych przedsięwzięć programistycznych.

% ============================================================

\end{document}
